\section{Every Gamma coefficient and an analytic normal generator}
\label{ung:section}

The canonical projection of Section~\ref{ibp:section} turns integration
by parts into a finite rational algorithm. The two universal-normal-form
continuations \cite{ung:reports} replace that algorithm, for the inverse-Gamma
hierarchy, by one beta integral for every coefficient. Their independent
proof routes give the same formula. We reconcile the rational-kernel route
here, then use its ordinary finite projections to obtain an analytic
generating identity for arbitrary holomorphic input jets. The latter also
specifies a domain on which the infinite derivative hierarchy converges.

\subsection{The universal coefficient identity}

Write $\operatorname{nf}$ for the projection of
Theorem~\ref{ibp:normal}, extended at weights zero and one as there.
In its differential polynomial ring put
\[
 \mathsf A_0=1,\qquad \mathsf A_{m+1}=u_0\mathsf A_m-\mathcal D\mathsf A_m,
 \qquad \mathsf N_m=\operatorname{nf}(\mathsf A_m).
\]
Its raw exponential generator and its projected finite jets are
\begin{align}
 \sum_{m\ge0}\mathsf A_m\frac{t^m}{m!}
 &=\exp\left(\sum_{r\ge0}\frac{(-1)^r u_rt^{r+1}}{(r+1)!}\right),
 \label{ung:raw}\\
 \sigma_d\sum_{m\ge0}\mathsf N_m\frac{t^m}{m!}
 &=e^{u_0t}F_d(X_1,\ldots,X_d),\qquad X_r=u_rt^{r+1},
 \label{ung:finite}
\end{align}
where $\sigma_d$ sets $u_r=0$ for $r>d$. The Appell factor follows
because translation in $u_0$ commutes with $\mathcal D$ and preserves
the canonical monomial space. Equations here initially mean formal
coefficient identities; analytic convergence is established below.

For $r\ge0$ define the explicit positive-coefficient polynomial
\begin{equation}\label{ung:C}
 C_r(v)=\frac{(v+1)^{r+1}-(v-1)^{r+1}}{2^{r+1}(r+1)!}
 =\sum_{\substack{0\le j\le r\\j\equiv r\ (2)}}
                   \frac{v^j}{2^rj!(r+1-j)!}.
\end{equation}
Thus $C_0=1$, $C_1=v/2$, $C_2=(1+3v^2)/24$,
and $2C_r'=C_{r-1}$.

\begin{theorem}[Every canonical inverse-Gamma coefficient]
\label{ung:coefficient}
The coefficient of $u_0^m$ in $\mathsf N_m$ is one. A monomial with
a unique largest positive derivative index is absent. For a canonical
monomial $u_0^{a_0}\cdots u_d^{a_d}$ with $d\ge1$, $a_d\ge2$ and
$m=\sum_{r=0}^d(r+1)a_r$, its coefficient vanishes when
$\sum r a_r$ is odd. Otherwise it is the finite rational number
\begin{equation}\label{ung:beta}
 \boxed{\frac{(-1)^d m!}{4^d d!(d+1)!\prod_{r=0}^d a_r!}
 \int_0^1(1-v^2)^d C_d(v)^{a_d-2}
                         \prod_{r=1}^{d-1}C_r(v)^{a_r}\,dv.}
\end{equation}
Every permitted coefficient is nonzero and has sign $(-1)^d$.
Consequently this gives the complete support, including all zeros,
at every weight and highest derivative order.
\end{theorem}

\begin{proof}
Use the Euler variational operator
$\mathcal E(P)=\sum_{r\ge0}(-\mathcal D)^r\partial_{u_r}P$.
Its commutator identity gives $\mathcal E\mathcal D=0$.
For a canonical polynomial, a monomial removed by $\sigma_d$ has
its highest index greater than $d$ at least twice. A partial derivative
removes at most one such factor, and subsequent total derivatives
leave a factor of index greater than $d$. Thus projection commutes with
the variational operation on canonical representatives, with
$\delta_d=\sum_{r<d}u_{r+1}\partial_{u_r}$ replacing $\mathcal D$.
This is why normal projection precedes jet truncation.

Applying $\mathcal E$ to \eqref{ung:raw} and integrating the formal
translation operator $e^{s\mathcal D}$ for $0\le s\le t$ gives,
after $s=t(1+v)/2$, the equation
\begin{equation}\label{ung:Euler-equation}
 F_d+\sum_{r=1}^d(-1)^r L_d^r\partial_{X_r}F_d
 =\frac12\int_{-1}^1e^{\sum_{r=1}^d C_r(v)X_r}\,dv,
 \qquad L_d=X_1+\sum_{r<d}X_{r+1}\partial_{X_r}.
\end{equation}
Here the conjugation by $e^{u_0t}$ turns $\delta_d$ into $t^{-1}L_d$.
Every operation is coefficientwise finite at a fixed weight.

This equation has at most one solution with canonical support. At
$d=1$ the homogeneous equation is $H-X_1H'=0$; its only possible
nonzero solution is linear, excluded by canonical support. Inductively,
specialization at $X_d=0$ reduces the homogeneous equation to the
lower-dimensional one, because the coefficient of $X_d$ is absent.
If $H$ is nonzero, let $c\ge2$ be its least occurring power of $X_d$.
The coefficient of $X_d^{c-1}$ is
$(-1)^d cL_{d-1}^dH_c=0$.
The operator $L_s$ is injective on formal series: for $s=1$ it is
multiplication by $X_1$, and its least power of $X_s$ has coefficient
killed by $L_{s-1}$ in the inductive case. Hence $H_c=0$, a contradiction.

For every $d\ge1$ put
\[
 \Delta_d(v)=(v+1)^{d+1}-(v-1)^{d+1},\qquad
 R_d(v)=\frac{(v-1)^{d+1}}{\Delta_d(v)},\qquad W_d(v)=2R_d'(v).
\]
Direct differentiation gives
\begin{equation}\label{ung:W}
 W_d(v)=\frac{4(d+1)(v^2-1)^d}{\Delta_d(v)^2}
       =\frac{(-1)^d(1-v^2)^d}{4^d d!(d+1)!C_d(v)^2}.
\end{equation}
The product rule telescopes to
\[
 \sum_{r=0}^d\frac{D_v^r((v-a)^{r+1}R')}{(r+1)!}
 =\frac{D_v^{d+1}((v-a)^{d+1}R)}{(d+1)!}-R.
\]
Subtract the identity at $a=1$ from the one at $a=-1$, using
$R=R_d$. Since $\Delta_dR_d=(v-1)^{d+1}$, this proves
\begin{equation}\label{ung:adjoint}
                    \sum_{r=0}^d2^rD_v^r(C_rW_d)=1.
\end{equation}

Let $J_d=\frac12\operatorname{FP}_v\int_{-1}^1
W_d(v)e^{\sum_{r=1}^dC_r(v)X_r}\,dv$.
For even $d$ the integral is ordinary. For odd $d$ its sole real
pole is at zero, with $W_d=-1/((d+1)v^2)+O(1)$; use the symmetric
cutoff and add $e^{\sum C_r(0)X_r}/((d+1)\epsilon)$ before taking
$\epsilon\downarrow0$. This fixes the finite part and permits
coefficientwise parameter differentiation.
The identity $L_d^re^{\sum C_rX_r}=2^rD_v^re^{\sum C_rX_r}$ and
\eqref{ung:adjoint} show by integration by parts that $J_d$ satisfies
\eqref{ung:Euler-equation}. The external boundary terms vanish because
$W_d$ has endpoint zeros of order $d$. Internal boundary terms have
zero constant coefficient: for a meromorphic boundary expression,
$B(-\epsilon)-B(\epsilon)$ has no constant term.

Normalize $e^{u_0t}J_d$ coefficientwise. The canonical projection
preserves its variational derivative, its Appell structure and its
largest possible index, so its normal part solves
\eqref{ung:Euler-equation} with canonical support. Uniqueness identifies
that part with \eqref{ung:finite}. Normalizing a term with a single
$u_d$ eliminates index $d$; terms with lower index cannot create it.
Thus coefficients containing $u_d^2$ are unchanged by normalization.
Their two factors of $C_d$ cancel the denominator in \eqref{ung:W},
leaving an ordinary polynomial integral. Its parity is $\sum r a_r$,
which proves \eqref{ung:beta}. At permitted parity all coefficients
of its remaining polynomial are nonnegative, with at least one positive,
so the integral is strictly positive. Finally the pure $u_0$ coefficient
is one, either from \eqref{ung:raw} or the constant term of
\eqref{ung:Euler-equation}.
\end{proof}

For example, expanding the polynomial integrand reduces every
coefficient to the elementary rational beta values
\[
 \int_0^1v^e(1-v^2)^d\,dv
 =\frac{2^dd!}{\prod_{j=0}^d(e+1+2j)}.
\]
The formula includes the mixed $u_0,u_1,u_2,u_3$ sectors and all
higher multiplicities without a new recurrence for each sector.

\subsection{Ordinary projections and analytic resummation}

Define the entire function
\[
 E_2(z)=\frac{e^z-1-z}{z^2},\qquad E_2(0)=\tfrac12,
 \qquad \kappa_r=\frac{(-1)^r}{4^r r!(r+1)!}.
\]
The disjoint highest-index sectors of Theorem~\ref{ung:coefficient}
give the ordinary entire finite projection
\begin{equation}\label{ung:ordinary-projection}
 F_d(\mathbf X)=1+\frac12\sum_{r=1}^d\kappa_rX_r^2
 \int_{-1}^1(1-v^2)^r
       e^{\sum_{j=1}^{r-1}C_j(v)X_j}E_2(C_r(v)X_r)\,dv.
\end{equation}
There are no poles on the integration interval. Its locally uniform
entire Taylor expansion gives exactly the coefficient formula above.

\begin{theorem}[An analytic normal generator for holomorphic jets]
\label{ung:analytic}
Let $U$ be holomorphic on $|z-x|<R$. Then
\[
 \mathsf G_U(t,x)=\sum_{m\ge0}
       \mathsf N_m\bigl(U(x),U'(x),\ldots\bigr)\frac{t^m}{m!}
\]
converges absolutely for $|t|<R$. It is the ordinary integral series
\begin{equation}\label{ung:analytic-generator}
 \boxed{\displaystyle
 e^{tU(x)}\left[1+\frac12\sum_{r\ge1}\kappa_rX_r^2
 \int_{-1}^1(1-v^2)^r
       e^{\sum_{j<r}C_j(v)X_j}E_2(C_r(v)X_r)\,dv\right],
 \quad X_r=U^{(r)}(x)t^{r+1}.}
\end{equation}
The sum is locally normally convergent on the indicated domain.
For varying $x$ it is jointly holomorphic wherever the corresponding
closed disks lie locally inside the domain of $U$.
If $U$ is entire, the generator is entire in $t$ at every $x$.
\end{theorem}

\begin{proof}
Fix $|t|<\rho<R$, let $M=\sup_{|z-x|\le\rho}|U(z)|$ and put
$q=|t|/\rho<1$. Cauchy's estimate gives
$|U^{(r)}(x)|\le Mr!\rho^{-r}$.
For real $|v|\le1$ the positive-coefficient polynomial in
\eqref{ung:C} satisfies $|C_r(v)|\le C_r(1)=1/(r+1)!$.
Consequently
\[
 |C_r(v)X_r|\le\frac{M|t|q^r}{r+1},\qquad
 \sum_{r\ge1}|C_r(v)X_r|\le\frac{M|t|q}{1-q}.
\]
The entire expansion $E_2(z)=\sum_{k\ge0}z^k/(k+2)!$ implies
$|E_2(z)|\le\tfrac12e^{|z|}$.
Using $\int_{-1}^1(1-v^2)^r\,dv\le2$, the $r$th summand in
\eqref{ung:analytic-generator}, including its external exponential,
is bounded in modulus by
\begin{equation}\label{ung:majorant}
 \frac{M^2|t|^2}{2(r+1)}
       e^{M|t|/(1-q)}\left(\frac{q^2}{4}\right)^r.
\end{equation}
This is a summable geometric majorant, uniform on every smaller
$t$ disk. Taking absolute Taylor coefficients before integration
gives the same majorant: replace all jet coefficients by their
Cauchy bounds and use the positive coefficients of $C_r$ and $E_2$.
Hence the double sum over highest indices and weights is absolutely
convergent. It can be regrouped by weight, and
Theorem~\ref{ung:coefficient} identifies those weights with
$\mathsf N_m/m!$. This proves the actual generating identity, rather
than only coefficientwise formal agreement.

Every integral summand is holomorphic, and the uniform majorant
permits termwise differentiation on smaller domains. Finite disk
covers give the assertion for varying $x$. For entire $U$ choose
$\rho>|t|$ arbitrarily. All sums at $t=0$ have their removable values.
\end{proof}

In particular $U=\psi$ is holomorphic on every disk avoiding
$0,-1,-2,\ldots$. Thus the canonical polygamma hierarchy has the
ordinary generator \eqref{ung:analytic-generator} for
\[
 |t|<\operatorname{dist}(x,\{0,-1,-2,\ldots\});
 \qquad x>0\ \Longrightarrow\ |t|<x.
\]
The normal projection remains the one in the previously established
differential identities; their primitive and finite-part formulas
retain the total-derivative endpoint terms.

For an affine input $U(x)=a+bx$ the integral series has the explicit
entire evaluation
\begin{equation}\label{ung:affine}
 \boxed{\mathsf G_U(t,x)=\frac{e^{t(a+bx)}}2
   \left(\cosh z+\frac{\sinh z}{z}-z\operatorname{Shi}(z)\right),
 \qquad z=\tfrac12bt^2,}
\end{equation}
where $\operatorname{Shi}(z)=\int_0^z\sinh(w)/w\,dw$ and
$\sinh z/z$ takes its removable value at zero. Only sector $r=1$
remains. With $W_1=(1-v^{-2})/2$, integration by parts in its fixed
symmetric finite part gives
\[
 F_1(2z)=\tfrac12\bigl(\sinh z/z+\cosh z-z\operatorname{Shi}(z)\bigr).
\]
The ordinary pole-free form \eqref{ung:ordinary-projection} gives the
same identity without a singular numerical quadrature. This supplies
all normalized affine Bell polynomials as coefficients of one
classical entire-function expression.
Equivalently, at every degree,
\begin{equation}\label{ung:affine-coefficients}
 \mathsf N_m(U,b,0,\ldots)=U^m-m!\sum_{k=1}^{\lfloor m/4\rfloor}
 \frac{U^{m-4k}b^{2k}}
 {(m-4k)!4^k(2k-1)(2k+1)(2k)!}.
\end{equation}
This follows by expanding the three entire functions in
\eqref{ung:affine}; it retains the canonical projection of the raw
affine Bell, or Hermite, hierarchy.

The bound also supplies a reproducible truncation prescription.
With $s=q^2/4$, the sum of sectors beyond $r=d$ is at most
\[
 \frac{M^2|t|^2e^{M|t|/(1-q)}}{2(d+2)(1-s)}s^{d+1}.
\]
For experimental evaluation use the pole-free expression
\eqref{ung:ordinary-projection} and the entire local series for $E_2$.
An independent sparse integration-by-parts implementation agrees with
the beta formula in all $1596$ tested coefficients through weight
eighteen. Twenty separate $75$-digit diagnostics compare ordinary
projections with independent normalized Taylor polynomials and the
truncation prescription, for real and complex Gamma inputs, an
exponential, a quadratic and a simple pole. The finite checks diagnose
the formulas; the all-order coefficient proof and the convergence
argument establish the infinite statements. No interval or
proof-assistant certification is asserted.
